\section{Background and Importance}
Professional training is essential for maintaining and improving employee competencies within an organization. It enables individuals to acquire new knowledge, adapt to technological changes, and take on greater responsibilities. For employees, it serves as a vehicle for career advancement and personal growth. For organizations, continuous training leads to improved productivity, better service quality, and stronger employee engagement, ultimately contributing to a competitive advantage in a rapidly evolving market.

\section{Problem Statement}
Despite its importance, many organizations still manage training activities through traditional, manual methods. The reliance on paper-based forms and static Excel spreadsheets introduces significant limitations in scalability and data management. These manual workflows are prone to human error, data duplication, and loss of information, making it extremely difficult to monitor ongoing activities and generate reliable performance indicators.

\section{Objectives}
The objective of this project is to design and implement a centralized web-based application, titled "TrainingMS", that automates the management of training activities. The system aims to:
\begin{itemize}
    \item Provide a secure platform for data entry and validation.
    \item Centralize participant, trainer, and session information.
    \item Simplify the scheduling and planning of training programs.
    \item Offer analytical tools to visualize and evaluate training outcomes through a statistics dashboard.
\end{itemize}
By digitizing these processes, the project seeks to improve operational efficiency, ensure data integrity, and support strategic decision-making for organizational stakeholders.
